\chapter{Exploration du Cadre Stratégique}

\section{Introduction}
Dans ce deuxième chapitre, nous approfondissons le cadre stratégique de notre projet...

\section{Contexte du projet}
Notre projet se situe dans le contexte des Objectifs de Développement Durable (ODD). Face au défi d'engager les utilisateurs, notre motivation est d'explorer des approches novatrices...

\section{Benchmarking}
Dans le cadre de l'exploitation des données, l'une des tâches clés consiste à comptabiliser avec précision les événements sur le réseau. Ce benchmark compare les algorithmes ML (SVM, KNN, Decision Tree, Random Forest) et Deep Learning (YOLOv8, Faster R-CNN, SSD).

\begin{figure}[H]
    \centering
    \includegraphics[width=0.7\textwidth]{media/f1_scores.png}
    \caption{Comparaison des F1-scores pour les différents algorithmes}
\end{figure}

Les modèles Deep Learning, notamment la famille YOLO, affichent d'excellentes performances. \textbf{YOLOv8s} obtient le meilleur score MOTA (0.300), tandis que les modèles classiques sont très rapides mais moins robustes sur le suivi d'objets complexes.

\begin{figure}[H]
    \centering
    \includegraphics[width=0.7\textwidth]{media/radar.png}
    \caption{Radar des KPIs globaux}
\end{figure}

\section{Discussion avec un inspecteur / professionnels}
% Insérez ici le texte de votre Word concernant la section 3 (Discussion)

\section{Business model canvas}
% Insérez ici le texte et l'image du BMC
% Exemple d'image :
% \begin{figure}[H]
%    \centering
%    \includegraphics[width=0.9\textwidth]{media/bmc.png}
%    \caption{Business Model Canvas}
% \end{figure}

\section{Conclusion}
% Insérez la conclusion du chapitre 2